In recent years quantum computing has witnessed a rapid increase in popularity, due to new solutions that can 
solve some of the most challenging problems, not only in computer science but also in many other fields~\cite{biondi2021quantum}.
However, a quantum computer with a sufficient amount of qubits is a serious threat to most of the current cryptographic
schemes, in particular asymmetric ones: quantum algoriths, such as Shor's~\cite{DBLP:conf/focs/Shor94}, can break the problems underlying
most of the current public-key cryptosystems in polynomial time, instead of the exponential time required by classical computers.\\
Since 2017 many institutions have started promoting the analysis of quantum-secure cryptosystems, with the most
prominent example being the Post-Quantum Cryptography (PQC) standardization contest organized by the U.S. National
Institute of Standards and Technologies (NIST). The standardization process started in 2016 and saw 
69 candidates from which, after three evaluation rounds, two standard candidates and four alternatives have emerged; 
the current standardization-elected candidate Key Encapsulation Mechanism (KEM) is Crystals-Kyber~\cite{DBLP:journals/iacr/BosDKLLSSS17}, 
a lattice-based cryptosystem.
To obtain resilience in the event of a potentially disruptive attack on lattice-based cryptosystems, the NIST extended the competition
with a fourth round, selecting four other candidates that rely on different computationally-hard problems; excluding SIKE, the remaining candidates,
namely BIKE~\cite{BIKE}, Classic McEliece~\cite{ClassicMcEliece} and HQC~\cite{melchor2018hamming}, are all based on 
the hardness of the syndrome decoding problem.

Hamming Quasi-Cyclic (HQC) is a code-based IND-CCA2 secure KEM whose security relies on the syndrome decoding problem for
quasi-cyclic codes. The scheme comprises both a Public Key Encryption scheme (PKE) and a Key Encapsulation Mechanism, 
obtained by applying the HHK Fujisaki-Okamoto transformation on the former. HQC uses a tensor code with a public generator
matrix $\mathbf{G}$ and a random double-circulant code with a systematic parity-check matrix; it is worth noting that the 
code generated by $\mathbf{G}$ is used only for its error-correcting capability, hence it is public and plays no role
in the security proofs.

\subsubsection{Our Contribution}
In this thesis we present the first fully constant-time software implementation of HQC, together with a comparative analysis
constant-time binary polynomial multiplication techniques, tailored for Cortex-M4 microprocessors. We provide an experimental
validation of the constant time properties via a battery of statistical t-tests, while varying the CPU frequency to highlight
the stability of our solution across different memory access latencies. Our results show that the optimal choice for constant-time
multiplications on the Cortex-M4 is a FFT-based approach, in contrast with classic divide-et-impera techniques (\textit{i.e.} Karatsuba)
which provide a better performance on x86 architectures.\\
Moreover, to counter the power attacks, we propose a masking scheme that can be applied to the entire cryptosystem, 
and present our findings on its potential efficacy; although our empirical results are inconclusive, they offer valuable insights
into our testing methodologies and pave the way for future research on the topic.

\subsubsection*{Organization}
The remaining of the thesis is organized as follows. In Chapter \ref{chap:background} we provide a brief introduction to the
mathematical concepts required to understand the multiplication techniques, side-channel attacks and the inner workings of HQC.
In Chapter \ref{chap:implementation} we describe our implementation of the multiplication algorithms and the masking scheme.
In Chapter \ref{chap:empirical} we present the experimental results of our implementation, and in Chapter \ref{chap:conclusions}
we draw the conclusions of our work.



